\begin{helper}
Fix the actual split-outside-blocker witness surface, a cut \(k<n\), a cell
\(\omega\), and a selected candidate \(x\in X\) with \(x\notin B_k\).
Let
\[
  (U^-_{k,\omega,y},U^+_{k,\omega,y})_{y\in X}
\]
be the first-witness residual pieces defined in
\(\noderef{SamplingSplitOutsideBlockerActualFirstWitnessResidualPieces}\).
Assume that a finite residual-label type \(\mathcal L\) has already been
chosen with paired source and target label events
\((E^-_\ell,E^+_\ell)_{\ell\in\mathcal L}\).  Assume that these labels
satisfy the concrete residual-label structure of
\(\noderef{SamplingSplitOutsideBlockerOutsideBkResidualLabelStructure}\),
that every source and target label event is null-measurable for the
corresponding product law, and that the exact pieces
\[
  U^-_{k,\omega,x}\cap C_\omega,
  \qquad
  U^+_{k,\omega,x}\cap C'_\omega
\]
are generated by the simultaneous truth vectors of these paired labels.
More precisely, there is then a finite index set \(\mathcal A\) and paired atoms
\[
  A^-_{x,a}\subseteq U^-_{k,\omega,x}\cap C_\omega,\qquad
  A^+_{x,a}\subseteq U^+_{k,\omega,x}\cap C'_\omega
\]
such that the two displayed pieces are the unions of their respective atom
families, the source atoms are pairwise almost-disjoint and null-measurable,
and the target atoms are pairwise almost-disjoint and null-measurable.

Each atom index \(a\) is a finite truth vector on \(\mathcal L\), and it
carries the truth function \(\tau_a:\mathcal L\to\{0,1\}\) such that
\[
  A^-_{x,a}
  =
  \{\eta:\eta\in E^-_\ell \text{ iff } \tau_a(\ell)=1
        \text{ for every }\ell\in\mathcal L\},
\]
and analogously
\[
  A^+_{x,a}
  =
  \{\eta':\eta'\in E^+_\ell \text{ iff } \tau_a(\ell)=1
        \text{ for every }\ell\in\mathcal L\}.
\]
The two residual pieces are recovered as the unions of those generated atoms,
so the residual piece is not an additional uninterpreted factor in the atom
normal form.

The finite label set is exhausted by the subfamilies deciding the selected
witness surface, every earlier first-witness exclusion, the cell predicates,
all retained noncurrent common-source blocker tests, all already-private
target blocker tests, and the global source and target boundary
complements.  The cited residual-label structure also gives coverage maps:
there are labels for the activation of the selected \(x\), for every
embedded-neighbor priority clause of \(x\), for every selected retained
source blocker test, for every earlier residual exclusion \(y<x\), for the
cell predicates, for every retained common source test \(k<i<n\), for every
already-private target test \(i<k\), and for the global boundary complements.
These labels are tied to the concrete source and target events of
\(\noderef{SamplingSplitOutsideBlockerActualFirstWitnessResidualPieces}\).
\end{helper}
\begin{proof}
Let the atom index set be the finite set of truth vectors
\(\mathcal A=\mathcal P_{\mathrm{fin}}(\mathcal L)\).  For a truth vector
\(a\subseteq\mathcal L\), define
\[
  A^-_{x,a}
  =
  \{\eta:\eta\in E^-_\ell\Longleftrightarrow \ell\in a
       \text{ for every }\ell\in\mathcal L\},
\]
and define \(A^+_{x,a}\) by the same formula with \(E^+_\ell\).  The truth
function attached to \(a\) is the Boolean indicator of membership in \(a\).
With these definitions, the atom normal forms required by
\(\noderef{SamplingSplitOutsideBlockerOutsideBkResidualLabelStructure}\) are
exactly the normal forms in the statement, so the assumed residual-label
structure applies to this atom family.

The two cover equalities are precisely the generation hypotheses: the source
piece \(U^-_{k,\omega,x}\cap C_\omega\) is the union of all source
truth-vector atoms, and the target piece \(U^+_{k,\omega,x}\cap C'_\omega\)
is the union of all target truth-vector atoms.  These hypotheses are the
place where the concrete selected-witness clauses, earlier exclusions, cell
predicates, retained common tests, already-private tests, and global
boundary complements are tied to the actual residual pieces.

It remains to prove pairwise almost-disjointness and null-measurability.  On
the source side, apply
\(\noderef{SamplingFiniteBooleanAtomPartition}\) with base event the whole
source product space and with label events \(E^-_\ell\).  The whole space is
null-measurable, and null-measurability of every \(E^-_\ell\) is one of the
hypotheses.  The cited Boolean partition lemma therefore says that the
truth-vector atoms for these source labels are pairwise almost-disjoint and
null-measurable.  Since intersecting with the whole space changes no atom,
these are exactly the source atoms \(A^-_{x,a}\).

The target side is identical.  Apply
\(\noderef{SamplingFiniteBooleanAtomPartition}\) to the whole target product
space and the target label events \(E^+_\ell\).  The target label
null-measurability hypotheses give null-measurability of each target atom,
and the Boolean partition lemma gives pairwise almost-disjointness.  Thus
the finite truth-vector atom family has the asserted source and target cover
equalities, paired residual-label structure, pairwise almost-disjointness,
and null-measurability.
\end{proof}
